\section{The Hodge ideals of simple normal crossing divisors}\label{SNC}
In this section we show that the Hodge ideals of divisors with simple normal crossing support essentially depend only on the support of the divisor, and therefore can be computed as in \cite[\S8]{MP1}.

\begin{proposition}\label{Hodge_ideals_SNC}
Let $X$ be a smooth variety, and $D$ an effective divisor on $X$ with simple normal crossing  support $Z$. Then for all $k$ we have 
$$I_k(D)=I_k(Z)\otimes_{\jepPubScript{O}_X} \jepPubScript{O}_X(Z - \lceil D \rceil).$$
\end{proposition}

\begin{proof}
Equivalently, we need to show that $I_k^\prime (D) = I_k(Z)$ for all $k$. The assertion is local, hence we may assume that we have coordinates $x_1,\ldots,x_n$
on $X$ such that $Z=H_1+\cdots+H_r$, where $H_i$ is defined by $x_i=0$. The morphism $X\to \jepPubBold{C}^r$ given by 
$(x_1,\ldots,x_r)$ is smooth, hence
using Remark~\ref{rem_behavior_etale} we see that it is enough to prove the proposition when 
$X= \operatorname{Spec}\,\jepPubBold{C}[x_1,\ldots,x_n]$ and $D=\sum_{i=1}^n\alpha_iH_i$,
where $H_i= {\rm div}(x_i)$ and $\alpha_i>0$. Let $\ell$ be the smallest positive integer such that all $a_i:=\ell\alpha_i$ are integers. The assertion to be proved is trivial 
when $\ell=1$, hence from now on we assume $\ell\geqslant 2$. 
Consider the Cartesian diagram 
$$
{\def\labelstyle{\textstyle}\xymatrix{V\ar[r]^{j_0}\ar[d]_{p}&W\ar[d]^{g}\\U\ar[r]_{j}&X,}}
$$
where 
$$j_0\colon V=\operatorname{Spec}\,\jepPubBold{C}[x_1^{\pm 1},\ldots,x_n^{\pm 1},y]/(y^{\ell}-x_1^{-a_1}\cdots x_n^{-a_n})$$
$$\longrightarrow W=\operatorname{Spec}\,\jepPubBold{C}[x_1,\ldots,x_n,z]/(z^{\ell}-x_1^{a_1}\cdots x_n^{a_n}),$$
with $j_0^*(z)=y^{-1}$. 
We will make use of some standard facts about cyclic covers with respect to simple normal crossing divisors,
exploiting the toric variety structure on the normalization of $W$. For basic facts regarding toric varieties, 
we refer to 
\cite{Fulton}.



Let $N$ be the lattice $\jepPubBold{Z}^n$ and $M$ its dual. We also consider the lattice 
$$N'=\{(v_1,\ldots,v_{n+1})\in\jepPubBold{Z}^{n+1}\mid a_1v_1+\cdots+a_nv_n=\ell v_{n+1}\}$$
and its dual 
$$M'=\jepPubBold{Z}^{n+1}/\jepPubBold{Z}\cdot (a_1,\ldots,a_n,-\ell).$$
Note that we have an injective lattice map $N'\to N$, with finite cokernel, induced by the projection onto the first $n$ components, and the dual map $M\to M'$
is again injective, with finite cokernel. In fact, we have an isomorphism $M'/M\simeq\jepPubBold{Z}/ \ell\jepPubBold{Z}$ that maps the class of $(u_1,\ldots,u_{n+1})\in M'$
to the class of $u_{n+1}$ in $\jepPubBold{Z}/\ell \jepPubBold{Z}$. 

We thus have an isomorphism $N'_{\jepPubBold{R}}\simeq N_{\jepPubBold{R}}=\jepPubBold{R}^n$. The strongly convex cone $\sigma=\jepPubBold{R}_{\geqslant 0}^n$ in 
$N_{\jepPubBold{R}}=\jepPubBold{R}^n$ gives the toric variety $X=\jepPubBold{C}^n$. As a cone in $N'_{\jepPubBold{R}}$, 
 $\sigma$ gives an affine toric variety $\widetilde{W}$, and the lattice map $N'\to N$ corresponds to a toric map
$\psi\colon \widetilde{W}\to X$. Note that we have a morphism of $\jepPubScript{O}(X)$-algebras 
$\jepPubScript{O}(W)\to \jepPubScript{O}(\widetilde{W})$ that maps $x_i$ to the element of $\jepPubBold{C}[\sigma^{\vee}\cap M']$
corresponding to the class of the $i$-th element of the standard basis of $\jepPubBold{Z}^n$, and $z$ to the class of $(0,\ldots,0,1)$. It is easy to check that
if we denote by $\lfloor \gamma\rfloor$ the largest integer $\leqslant\gamma$, then
\begin{equation}\label{eq_decomp_W}
\jepPubScript{O}(\widetilde{W})=\mathop{\textstyle\bigoplus}\limits_{0\leqslant j\leqslant\ell-1}\jepPubScript{O} (X)x_1^{-\lfloor j\alpha_1\rfloor}\cdots x_n^{-\lfloor j\alpha_n\rfloor}z^j,
\end{equation}
and consequently to deduce that $\jepPubScript{O}(\widetilde{W})$ is integral over $\jepPubScript{O}(W)$. As the coordinate ring of a toric variety, $\jepPubScript{O}(\widetilde{W})$
is normal, hence it is the integral closure of $\jepPubScript{O}(W)$ in its field of fractions.
Moreover, since $\widetilde{W}$ is a toric variety, we may choose a toric resolution of singularities $Y\to \widetilde{W}$, and let $f\colon Y\to X$ be the composition. Since the map
$Y\to W$ is an isomorphism over the complement of $g^{-1} (\sum H_i)$, it follows that there is an open embedding $\iota\colon V\hookrightarrow Y$ such that $f\circ\iota =j\circ p$. The support $E_Y$ of $Y\smallsetminus
\iota(V)$ is the sum of all prime toric divisors on $Y$. 

The action of the torus $T_{M'}={\rm Spec}\,\jepPubBold{C}[M']$ on $\widetilde{W}$ induces an action of the finite group
${\rm Spec}\,\jepPubBold{C}[M'/M]\simeq {\rm Spec}\,\jepPubBold{C}[\jepPubBold{Z}/\ell\jepPubBold{Z}]=\mu_{\ell}$ on 
$\widetilde{W}$. This is the action induced on the normalization $\widetilde{W}$ by the $\mu_{\ell}$-action on $W$ that we discussed in Section~\ref{coverings}. In particular,
the toric resolution $Y\to\widetilde{W}$ is automatically equivariant. Note that in the decomposition (\ref{eq_decomp_W}), an element $\lambda\in\mu_{\ell}$ acts on the 
summand corresponding to $j$ by multiplication with $\lambda^j$.


The equality $f\circ\iota =j\circ p$ implies that we have an isomorphism of filtered $\jepPubScript{D}_X$-modules
$$j_+p_+\jepPubScript{O}_V\simeq f_+\iota_+\jepPubScript{O}_V=f_+\jepPubScript{O}_Y(*E_Y).$$
As usual, in order to compute the push-forward of $\jepPubScript{O}_Y(*E_Y)$, it is more convenient to work with right $\jepPubScript{D}$-modules.
Recall that there is a complex of right $\jepPubScript{D}_Y$-modules
$$A^\jepPubBullet = A_Y^{\jepPubBullet}:\,\,\,0 \longrightarrow \jepPubScript{D}_Y \longrightarrow \Omega_Y^1(\log E_Y) \otimes_{\jepPubScript{O}_Y} \jepPubScript{D}_Y \longrightarrow \cdots  \longrightarrow \omega_Y(E_Y) \otimes_{\jepPubScript{O}_Y} \jepPubScript{D}_Y \longrightarrow 0$$
located in degrees $-n,\ldots,0$, that is filtered quasi-isomorphic to $\omega_Y(*E_Y)$; see the beginning of Section~\ref{complex_for_SNC}. 
Since $Y$ is a toric variety, we have a canonical isomorphism 
$\Omega_Y^1(\log E_Y)\simeq M'\otimes_{\jepPubBold{Z}}\jepPubScript{O}_Y$ (see \cite[\S4.3]{Fulton}). 
We will also consider the corresponding complex on $X$:
$$A_X^{\jepPubBullet}:\,\,\,0\longrightarrow\jepPubScript{D}_X\longrightarrow M\otimes_{\jepPubBold{Z}}\jepPubScript{D}_X\longrightarrow \cdots\longrightarrow \wedge^nM\otimes_{\jepPubBold{Z}}\jepPubScript{D}_X\longrightarrow 0.$$

It follows from the definition that,  forgetting about the filtration, we have
$$f_+\omega_Y(*E_Y)=\jepPubBold{R} f_*(A^{\jepPubBullet}\otimes_{\jepPubScript{D}_Y}\jepPubScript{D}_{Y\to X}).$$
Note that $\jepPubScript{D}_{Y\to X}=f^* \jepPubScript{D}_X$ as $\jepPubScript{O}_Y$-modules, hence the projection formula implies
$$R^if_*(A^{p-n}\otimes_{\jepPubScript{D}_Y}\jepPubScript{D}_{Y\to X})\simeq \wedge^pM'\otimes_{\jepPubBold{Z}}R^if_*\jepPubScript{O}_Y\otimes_{\jepPubScript{O}_X}\jepPubScript{D}_X=0$$
for $i>0$, since $f$ is the composition of a finite map with a toric resolution. Therefore 
$f_+\omega_Y(*E_Y)$ is represented by the complex $B^{\jepPubBullet}$, where
$$B^{p-n}=\wedge^pM'\otimes_{\jepPubBold{Z}}\psi_*\jepPubScript{O}_{\widetilde{W}}\otimes_{\jepPubScript{O}_X}\jepPubScript{D}_X.$$
In order to describe the differential of this complex, it is convenient to use the isomorphism $M_{\jepPubBold{Q}}\simeq M'_{\jepPubBold{Q}}$
and the decomposition (\ref{eq_decomp_W}). With a little care, it follows from the definitions that if we put
$$B^{p-n}_j=\wedge^pM_{\jepPubBold{Q}}\otimes_{\jepPubBold{Q}}\jepPubScript{O}_X\cdot x_1^{-\lfloor j\alpha_1\rfloor}\cdots x_n^{-\lfloor j\alpha_n\rfloor}z^j\otimes
_{\jepPubScript{O}_X}\jepPubScript{D}_X,$$
then $B^{\jepPubBullet}$ decomposes as the direct sum of the subcomplexes $B^{\jepPubBullet}_j$, for $j$ such that $0\leqslant j\leqslant \ell-1$. 
Furthermore, if we identify each $B^{p-n}_j$ in the obvious way with $A^{p-n}_X$, then the differential 
$$\delta_{B_j}^{p-n}\colon \wedge^pM_{\jepPubBold{Q}}\otimes_{\jepPubBold{Q}}\jepPubScript{D}_X\longrightarrow \wedge^{p+1} M_{\jepPubBold{Q}}\otimes_{\jepPubBold{Q}}\jepPubScript{D}_X$$
is given by 
$$\delta_{B_j}^{p-n}=\delta_{A_X}^{p-n}+(w_j \wedge -)\otimes {\rm Id}_{\jepPubScript{D}_X},$$
where $\delta_{A_X}$ is the differential on $A_X^{\jepPubBullet}$ and
$$w_j=(w_{j,1},\ldots,w_{j,n}),\quad\text{with}\quad w_{j,i}=j\alpha_i-\lfloor j\alpha_i\rfloor.$$
It follows from Proposition~\ref{prop_SNC_complex} that we have a morphism
$$B_j^0\longrightarrow \jepPubScript{M}_r(x_1^{w_{j,1}}\cdots x_n^{w_{j,n}})$$ 
that induces a quasi-isomorphism
$$B_j^{\jepPubBullet}\longrightarrow \jepPubScript{M}_r(x_1^{w_{j,1}}\cdots x_n^{w_{j,n}})$$
(see also Remark~\ref{twist_individual}).

We now bring the filtrations into the picture. It follows from Saito's strictness results (see the discussion in 
Section~\ref{Hodge_modules}; cf. also \cite[\S4,~\S6]{MP1}) that 
$$F_kf_+\omega_Y(*E_Y)={\rm Im}\big(\jepPubBold{R} f_*F_k(A^{\jepPubBullet}\otimes_{\jepPubScript{D}_Y}\jepPubScript{D}_{Y\to X})\longrightarrow \jepPubBold{R} f_*(A^{\jepPubBullet}\otimes_{\jepPubScript{D}_Y}\jepPubScript{D}_{Y\to X})\big).$$
Arguing as above, we deduce that 
$$F_kf_+\omega_Y(*E_Y)={\rm Im}\big(f_*F_k(A^{\jepPubBullet}\otimes_{\jepPubScript{D}_Y}\jepPubScript{D}_{Y\to X})\longrightarrow f_*(A^{\jepPubBullet}\otimes_{\jepPubScript{D}_Y}\jepPubScript{D}_{Y\to X})\big).$$
In other words, $(f_+\omega_Y(*E_Y), F)$ is represented by the filtered complex $B^{\jepPubBullet}$, and using Proposition~\ref{prop_SNC_complex},
we conclude that 
$$f_+\omega_Y(*E_Y)\simeq \mathop{\textstyle\bigoplus}\limits_{j=0}^{\ell-1}\jepPubScript{M}_r(x_1^{w_{j,1}}\cdots x_n^{w_{j,n}}),$$
where the filtration on $\jepPubScript{M}_r(x_1^{w_{j,1}}\cdots x_n^{w_{j,n}})$ is given by 
$$F_{-n}\jepPubScript{M}_r(x_1^{w_{j,1}}\cdots x_n^{w_{j,n}})=x_1^{w_{j,1}}\cdots x_n^{w_{j,n}}\omega_X(Z)\quad\text{and}$$
$$F_{k-n}\jepPubScript{M}_r(x_1^{w_{j,1}}\cdots x_n^{w_{j,n}})=F_{-n}\jepPubScript{M}_r(x_1^{w_{j,1}}\cdots x_n^{w_{j,n}})\cdot F_k\jepPubScript{D}_X\quad\text{for}\quad k\geqslant 1.$$

By comparing the $\mu_{\ell}$-actions, we conclude that the summand 
$\jepPubScript{M}_r(x_1^{-\alpha_1}\cdots x_n^{-\alpha_n})$ on which an element $\lambda\in\mu_{\ell}$ acts by multiplication with $\lambda^{-1}$
corresponds to $j=\ell-1$. Therefore the filtration on 
$\jepPubScript{M}(x_1^{-\alpha_1}\cdots x_n^{-\alpha_n})$ is given by
$$F_{-n}\jepPubScript{M}_r(x_1^{-\alpha_1}\cdots x_n^{-\alpha_n})=(x_1^{-\alpha_1}\cdots x_n^{-\alpha_n})x_1^{\lceil\alpha_1\rceil}\cdots x_n^{\lceil\alpha_n\rceil}\omega_X(Z)\quad\text{and}$$
$$F_{k-n}\jepPubScript{M}_r(x_1^{-\alpha_1}\cdots x_n^{-\alpha_n})=F_{-n}\jepPubScript{M}_r(x_1^{-\alpha_1}\cdots x_n^{-\alpha_n})\cdot F_k\jepPubScript{D}_X\quad\text{for}\quad k\geqslant 1.$$
It is now a straightforward computation to see that
$I_k^\prime(D)$ is the ideal generated by the monomials $\prod_{i=1}^nx_i^{c_i}$, where $0\leqslant c_i\leqslant k$ for all $i$ and $\sum_ic_i=(n-1)k$. 
This coincides with $I_k (Z)$ according to \cite[Prop.~8.2]{MP1},
completing the proof of the proposition.
\end{proof}







